\chapter{Overview of Popular Time-Series Databases}

This chapter provides an overview of the most widely used time-series
databases. For each database, the core architecture, query language, and key
features relevant to time-series workloads are described.

\section{InfluxDB}

InfluxDB is a stand-alone time-series database developed by InfluxData.
Version~1.0, released in 2016, was written in Go and used a custom storage
engine called the Time-Structured Merge Tree (TSM), optimised for metrics
workloads. InfluxDB 3 is a complete rebuild of the core database engine written
in Rust, using the FDAP stack (Apache Flight, DataFusion, Apache Arrow, and
Parquet). Rust was chosen for its superior performance, memory safety, and
fearless concurrency compared to the previous Go implementation. InfluxDB 3
Core is licensed under the MIT/Apache~2.0 license and a commercial Enterprise
edition is also available.~\cite{InfluxDB:01}

InfluxDB~3 removes cardinality limits, introduces tiered object storage, and
adds SQL as a first-class query language. It also includes the InfluxDB
Explorer, a UI for querying, exploring, and visualising data.

It can be queried with SQL or with the InfluxQL query
language.~\cite{InfluxDB:02}

\subsection{InfluxQL}

InfluxQL is the SQL-like query language that can be used to query InfluxDB\@.
Its syntax aims to feel familiar to SQL while offering specialised features for
time series data. Although SQL is now the primary query language in InfluxDB~3,
InfluxQL remains supported for backwards compatibility.

The following is an example query to list tag values for a specific key:

\begin{lstlisting}[language=SQL, caption={InfluxQL query to show tag values.}, label={lst:influxql-query}]
SHOW TAG VALUES FROM "measurement_table" WITH KEY = "tag-key" WHERE time > now() - 2d
\end{lstlisting}

\section{TimescaleDB}

TimescaleDB is an open-source~\cite{TimescaleDB:07} Postgres extension that
allows one to turn any Postgres database into a time series database. Its main
advantage is the Postgres core, which makes it suitable for applications that
need improved time series data performance while retaining a fully capable
relational database.

It supports standard SQL via its Postgres core, with a few extensions for
time-series use cases.~\cite{TimescaleDB:08}

\subsection{Time-series-specific SQL functions}

TimescaleDB extends standard SQL with time-series specific functions to make it
easier to work with time-series data. Here is an example query that aggregates
data over time intervals:

\begin{lstlisting}[language=SQL, caption={TimescaleDB query to aggregate data over time intervals.}, label={lst:timescaledb-query}]
SELECT time_bucket('1 hour', time) AS bucket,
    AVG(value) AS average_value
FROM measurement_table
WHERE time > now() - interval '1 day'
GROUP BY bucket
ORDER BY bucket;
\end{lstlisting}

\subsection{Hypertables}

The easiest way to optimize a Postgres table for time-series data in
TimescaleDB is to convert it into a hypertable. Hypertables improve performance
and storage efficiency by partitioning the data based on timestamps. For
details on how hypertables are set up and used in the AirScout project, see
Section~\ref{sec:timescaledb-features}.

~\cite{TimescaleDB:02}

\subsection{Hypercore}

The storage engine of Postgres is well optimized for general use cases, but not
specifically for time-series data. Its row-based architecture allows for fast
reads and writes of individual rows, but can be inefficient for time-series
data that is often queried in bulk over time ranges. Hypercore is Timescale's
hybrid storage engine built to address this by combining the best of both row-
and column-oriented storage.

Hypercore operates in two distinct layers:
\begin{itemize}
  \item \textbf{Row-oriented layer (hot data):} Recent data is stored in a standard row format, optimized for fast inserts and low-latency access to the most current readings.
  \item \textbf{Column-oriented layer (cold data):} Older data chunks are automatically transitioned into a compressed columnar format. This dramatically reduces storage footprint and speeds up analytical queries that scan many rows across a time range, since only the relevant columns need to be read from disk.
\end{itemize}

This hybrid approach allows TimescaleDB to deliver high performance for both
real-time ingestion and large-scale historical queries. Because compression
algorithms such as delta encoding and run-length encoding work best on
sequences of similar values found in a single column, the columnar format also
achieves significantly higher compression ratios than row-based storage.

Using Hypercore requires only enabling it on a hypertable:

\begin{lstlisting}[language=PGTSSQL, caption={Enabling Hypercore on a TimescaleDB hypertable.}, label={lst:timescaledb-hypercore}]
CREATE TABLE crypto_ticks (
   "time" TIMESTAMPTZ,
   symbol TEXT,
   price DOUBLE PRECISION,
   day_volume NUMERIC
) WITH (
  tsdb.hypertable,
  tsdb.partition_column='time',
  tsdb.segmentby='symbol', 
  tsdb.orderby='time DESC'
);
\end{lstlisting}
~\cite[TimescaleDB Hypercore]{TimescaleDB:04}

\subsection{How to use TimescaleDB}

There are multiple ways to install and use TimescaleDB\@. A very common way is
to use Docker.

Here is an example Docker Compose service definition for TimescaleDB\@:

\begin{lstlisting}[style=yaml, caption={Docker-compose service definition for TimescaleDB.}, label={lst:docker-timescaledb}]
services:
  timescaledb:
    image: timescale/timescaledb:latest-pg14
    container_name: timescaledb-dev
    environment:
      POSTGRES_DB: airscout
      POSTGRES_USER: airscout_user
      POSTGRES_PASSWORD: <password>
    ports:
      - "5434:5432"
    volumes:
      - timescaledb_data_dev:/var/lib/postgresql/data
      - ./AirscoutDB/migrations:/docker-entrypoint-initdb.d/
    networks:
      - airscout_network_dev
    healthcheck:
      test: ["CMD-SHELL", "pg_isready -U airscout_user -d airscout"]
      interval: 10s
      timeout: 5s
      retries: 5
volumes:
  timescaledb_data_dev:
networks:
  airscout_network_dev:
\end{lstlisting}

\section{Prometheus}

Prometheus is an open-source time series database, created by SoundCloud in
2012 to improve their monitoring capabilities. It is specialized for monitoring
and very popular together with Grafana, a popular open-source platform for
monitoring and observability. The database is written in Go and is licensed
under the Apache 2.0 license. In contrast to other time series databases,
Prometheus stores all data as time series data. It is not recommended to insert
logs manually, instead use tools like Grafana Loki or OpenSearch to feed logs
into the system.

~\cite{Prometheus:01}
~\cite{Prometheus:03}

\section{PromQL}

Prometheus can be queried using its own functional query language called PromQL
(Query Language of Prometheus). Unlike SQL, PromQL is specifically designed for
time-series data and works on the concept of metrics and labels. A metric is a
named measurement \\ (e.g., \texttt{http\_requests\_total}), and labels are
key-value pairs attached to each metric that allow filtering and grouping by
dimensions such as host, region, or status code. PromQL expressions return one
of four data types: instant vectors (a single value per matching time series at
a given moment), range vectors (a sequence of values over a time range),
scalars, or strings.\cite{Prometheus:02}

The following PromQL query shows the unused memory in MiB for each instance, as
reported by a fictional cluster scheduler exposing instance metrics:

\begin{lstlisting}[language=SQL, caption={PromQL query to get unused memory in MiB for every instance.}, label={lst:promql-query}]
(instance_memory_limit_bytes - instance_memory_usage_bytes) / 1024 / 1024
\end{lstlisting}

Explanations:

\begin{itemize}
  \item
        \texttt{instance\_memory\_limit\_bytes} and \texttt{instance\_memory\_usage\_bytes} are metric \\names that represent the total memory limit and the current memory usage of each instance, respectively. In Prometheus, metrics are time-series data identified by a name and optional key-value pairs called labels. Since both metrics share the same label set (e.g.,~\texttt{instance},~\texttt{job}), Prometheus automatically matches them element-wise.
  \item
        The subtraction operation \texttt{-} calculates the difference between the total memory limit and the current memory usage, resulting in the unused memory in bytes.
  \item
        Dividing by \texttt{1024 / 1024} converts the result from bytes to mebibytes (MiB), producing a human-readable output for monitoring dashboards.
\end{itemize}

\section{OpenTSDB}

Most of the previously discussed databases are designed to only run on a single
machine. OpenTSDB differs in this as it can run on multiple machines in a
distributed fashion. It is built on top of HBase, which is a distributed,
scalable, big data store that runs on top of the Hadoop Distributed File System
(HDFS). OpenTSDB is written in Java and is licensed under the LGPLv2.1 or the
GPL-3.0 license.

Data can be ingested using Telnet, the HTTP API, or a client library.
~\cite{OpenTSDB:01}

Example Query via HTTP API\@:

\begin{lstlisting}[language=SQL, caption={OpenTSDB HTTP API query to get average CPU usage over the last hour.}, label={lst:opentsdb-query}]
http://localhost:4242/api/query?start=1h-ago&m=avg:cpu.usage{host=server1}
\end{lstlisting}

\section{MongoDB}

MongoDB is a document-oriented NoSQL database that stores data in flexible,
JSON-like documents. Its architecture emphasizes application-oriented data
models, built-in replication for high availability, and horizontal scaling via
sharding. Although MongoDB is not a dedicated time-series database, it supports
time-series workloads through dedicated time series collections and is often
considered when flexible schemas or existing MongoDB deployments are already in
place.~\cite{MongoDB:01} The MongoDB server is primarily written in C++ and is
published under the SSPL.~\cite{MongoDB:02} The SSPL means the software is free
to use, modify, and distribute, but if you offer the software as a service, you
must release the source code of your service under the same license.

Interestingly, MongoDB Inc., the company behind MongoDB, generated 2.01 billion
USD in revenue in 2025, with a significant portion likely coming from their
cloud services that include MongoDB Atlas, which offers time-series data
management capabilities. This highlights the commercial success and widespread
adoption of MongoDB in various applications.~\cite{MongoDB:05}

MongoDB can be queried using standard document queries and the MongoDB Query
Language (MQL) aggregation pipeline. For time-series analysis, the aggregation
pipeline is particularly relevant because it supports filtering, grouping, and
time bucketing operations on measurement data.~\cite{MongoDB:01}

\subsection{Time Series Collections}

Since MongoDB~5.0, time series collections provide a storage model optimized
for time-stamped measurements. Each measurement contains a \texttt{timeField},
the observed metric values, and optionally a \texttt{metaField} that stores
metadata identifying the data source, such as a sensor ID or location.
Internally, MongoDB organizes this data in a time-oriented, columnar storage
format that improves query efficiency, reduces disk usage, and lowers I/O for
read-heavy workloads compared to regular collections.~\cite{MongoDB:03}

\subsection{Aggregation Pipeline}

The following example query groups measurements into one-hour buckets and
calculates the average value for each bucket:

\begin{lstlisting}[caption={MongoDB aggregation query to aggregate data over time intervals.}, label={lst:mongodb-query}]
db.measurement_table.aggregate([
  {
    $match: {
      timestamp: { $gt: new Date(Date.now() - 24 * 60 * 60 * 1000) }
    }
  },
  {
    $group: {
      _id: {
        bucket: {
          $dateTrunc: { date: "$timestamp", unit: "hour" }
        }
      },
      average_value: { $avg: "$value" }
    }
  },
  { $sort: { "_id.bucket": 1 } }
])
\end{lstlisting}

This query uses \texttt{\$dateTrunc} to round each timestamp down to the start
of the hour and then applies \texttt{\$avg} to compute the mean value per
bucket.~\cite{MongoDB:04}

MongoDB's main advantages for time-series workloads include its flexible
schema, rich query capabilities, and strong support for horizontal scaling.
However, it does not match the performance of specialized time-series databases
for high-cardinality or write-heavy workloads due to its general-purpose design
(see Chapter~\ref{sec:benchmarking}).